\section{Trace-class transfer and holomorphic de Rham coupling}
\label{sec:holomorphic}

\subsection{Marked atom transfer}

Fix \(\eta>\half\) and define
\begin{equation}
T_\eta(s,u)
=\sum_{p\in\At(P)}u^{\ell(p)}p^{-s}q_p.
\label{eq:atom-transfer}
\end{equation}
The summation index is derived by the cover predicate, not read from a prime
table.

\begin{theorem}[Holomorphic trace-class family]
\label{thm:trace-class-transfer}
Whenever
\begin{equation}
\sum_{p\in\At(P)}
|u|^{\ell(p)}p^{-\Re s}<\infty,
\label{eq:absolute-domain}
\end{equation}
the series \eqref{eq:atom-transfer} converges in trace norm, locally uniformly
in \((s,u)\). For every \(r\ge1\),
\begin{equation}
T_\eta(s,u)^r
=\sum_p u^{r\ell(p)}p^{-rs}q_p,
\qquad
\Tr T_\eta(s,u)^r
=\sum_p u^{r\ell(p)}p^{-rs}.
\label{eq:all-power-traces}
\end{equation}
Its Fredholm determinant is
\begin{equation}
\det(I-zT_\eta(s,u))
=\prod_p \left(1-zu^{\ell(p)}p^{-s}\right).
\label{eq:atom-fredholm}
\end{equation}
\end{theorem}

\begin{proof}
The uniform bound \eqref{eq:uniform-trace-norm} and
\eqref{eq:absolute-domain} prove absolute trace-norm convergence, locally
uniformly on compact subsets of the domain. Holomorphy as an
\(\Sone\)-valued map follows. Pair annihilation from
\cref{thm:rank-one} gives the operator identity in
\eqref{eq:all-power-traces}; tracing uses \(\Tr q_p=1\).

For small \(z\), the trace-class determinant identity
\[
\log\det(I-zT)
=-\sum_{r\ge1}\frac{z^r}{r}\Tr T^r
\]
and absolute convergence give
\[
\log\det(I-zT_\eta)
=\sum_p \log\left(1-zu^{\ell(p)}p^{-s}\right).
\]
Both sides of \eqref{eq:atom-fredholm} are entire in \(z\), so equality
extends globally.
\end{proof}

For \(\rho=|u|>0\), the gamma duration has the two-sided comparison
\begin{equation}
\rho^{\ell(n)}
\asymp_\rho n^{2\log_2\rho}.
\label{eq:gamma-asymptotic}
\end{equation}
The prime series therefore gives the exact threshold
\begin{equation}
\Re s>1+2\log_2\rho.
\label{eq:marker-half-plane}
\end{equation}
At \(u=0\) the transfer is identically zero. A common, marker-transparent safe
domain is \(|u|\le1\) and \(\Re s>1\).

At \(u=z=1\), \eqref{eq:atom-fredholm} becomes
\begin{equation}
\det(I-T_\eta(s,1))
=\prod_p(1-p^{-s})
=\frac1{\zeta(s)},
\qquad \Re s>1.
\label{eq:reciprocal-zeta}
\end{equation}
Equation \eqref{eq:reciprocal-zeta} is a consequence of the source-derived
loop theorem. It is not advertised as a new scalar identity.

\begin{proposition}[Sharp \(u=1\) trace-class barrier]
\label{prop:trace-barrier}
The realization \(T_\eta(s,1)\) cannot be trace class when \(\Re s\le1\).
\end{proposition}

\begin{proof}
For each source atom \(p\), the vector \(r_p=e_1+e_p\) satisfies
\[
T_\eta(s,1)r_p=p^{-s}r_p.
\]
These vectors are linearly independent. The eigenvalues of a trace-class
operator, with algebraic multiplicity, are absolutely summable. Their
absolute values here include \(p^{-\Re s}\), whose sum over primes diverges
for \(\Re s\le1\).
\end{proof}

The marker may shift the absolute half-plane when \(|u|<1\), but this does
not continue the \(u=1\) arithmetic operator. In particular, a scalar
continuation of the right side of \eqref{eq:reciprocal-zeta} is not a
same-object continuation of \(T_\eta(s,1)\).

\subsection{Two honest holomorphic degrees}

We now couple the incidence selector to the inherited prefix-coded affine
renewal sector. Let \(U_{p,0}\) and \(U_{p,1}\) denote its zero- and one-form
pullbacks on the fixed holomorphic cylinder space. The prior construction
provides a common degreewise trace-norm bound on the absolute domain and, for
every branch word \(w\), the local Lefschetz identity
\begin{equation}
\Tr U_{w,0}-\Tr U_{w,1}=1
\label{eq:local-lefschetz}
\end{equation}
after removing the scalar branch weight. Concretely, if the derivative of the
affine word is \(\vartheta_w\), the two local traces are
\[
\frac1{1-\vartheta_w},
\qquad
\frac{\vartheta_w}{1-\vartheta_w},
\]
whose difference is one. This is the canonical de Rham cancellation of the
fixed-point denominator.

Define the degreewise operators
\begin{equation}
\cT_k(s,u)
=\sum_p u^{\ell(p)}p^{-s}q_p\otimes U_{p,k},
\qquad k=0,1.
\label{eq:degreewise-transfer}
\end{equation}

\begin{theorem}[Honest graded determinant]
\label{thm:graded-determinant}
On the common absolute trace-class domain, both operators in
\eqref{eq:degreewise-transfer} are trace class and, for every \(r\ge1\),
\begin{equation}
\Tr\cT_0(s,u)^r-\Tr\cT_1(s,u)^r
=\sum_p u^{r\ell(p)}p^{-rs}.
\label{eq:graded-power-trace}
\end{equation}
Consequently
\begin{equation}
\Dgr(s,u,z)
:=\frac{\det(I-z\cT_0(s,u))}
        {\det(I-z\cT_1(s,u))}
=\prod_p \left(1-zu^{\ell(p)}p^{-s}\right).
\label{eq:graded-ratio}
\end{equation}
\end{theorem}

\begin{proof}
The common local trace-norm bound, \eqref{eq:absolute-domain}, and
\eqref{eq:uniform-trace-norm} give degreewise trace-norm convergence. In an
\(r\)-fold product, the incidence factor annihilates every mixed atom word
before the holomorphic trace is taken. A monochromatic word contributes
\(\Tr q_p=1\), and \eqref{eq:local-lefschetz} removes its local stability
factor. This proves \eqref{eq:graded-power-trace}. Taking the difference of
the two ordinary trace-log expansions yields \eqref{eq:graded-ratio}.
\end{proof}

\begin{figure}[t]
\centering
\resizebox{0.94\textwidth}{!}{\begin{tikzpicture}[
  font=\sffamily\small,
  box/.style={draw=charcoal,rounded corners=2pt,very thick,align=left,
    inner sep=7pt},
  pass/.style={box,fill=passgreen!9,draw=passgreen},
  warn/.style={box,fill=warningamber!10,draw=warningamber},
  stop/.style={box,fill=stopred!8,draw=stopred},
  lab/.style={font=\sffamily\scriptsize,text=charcoal},
  arr/.style={-{Latex[length=1.5mm,width=1.1mm]},very thick,draw=charcoal}
]
\node[font=\sffamily\bfseries] at (2.6,2.7) {weighted incidence fiber};
\node[pass,minimum width=64mm] (individual) at (2.6,1.55) {
  $\eta>\half$\\
  each $q_p$ is rank-one trace class\\
  $\norm{q_p}_1\le\sqrt{2C_\eta}$};
\node[pass,minimum width=64mm] (global) at (2.6,-0.55) {
  $\eta>1$\\
  $Z_\eta,M_\eta$ bounded inverses\\
  whole family similar to coordinates};

\node[font=\sffamily\bfseries] at (10.5,2.7) {marked transfer family};
\node[pass,minimum width=76mm] (domain) at (10.5,1.4) {
  $\Re s>1+2\log_2|u|$\\
  $T_\eta(s,u)$ and both de Rham degrees\\
  converge in trace norm};
\node[warn,minimum width=76mm] (safe) at (10.5,-0.6) {
  common safe domain\\
  $|u|\le1,\ \Re s>1$};
\node[stop,minimum width=76mm] (barrier) at (10.5,-2.45) {
  at $u=1$: stop at $\Re s=1$\\
  eigenvalues $p^{-s}$ are not absolutely summable};

\draw[arr] (individual.south) -- (global.north);
\draw[arr] (domain.south) -- (safe.north);
\draw[arr] (safe.south) -- (barrier.north);
\draw[very thick,draw=stopred,dashed] (6.65,-3.55) -- (14.35,-3.55);
\node[lab,text=stopred] at (10.5,-3.95)
  {scalar continuation of the product is not continuation of this operator};
\end{tikzpicture}
}
\caption{The theorem domains have different owners. Individual incidence
idempotents exist for \(\eta>\half\); a bounded global incidence similarity
is proved for \(\eta>1\); and the marked transfer plus both holomorphic
degrees are trace class in the stated \((s,u)\)-half-plane. At \(u=1\), the
eigenvalue obstruction stops the operator at \(\Re s=1\), irrespective of
any scalar continuation of its product.}
\label{fig:analytic-domains}
\end{figure}

\begin{remark}[Ordinary versus graded]
\label{rem:ordinary-graded}
The numerator and denominator in \eqref{eq:graded-ratio} are separately
defined ordinary Fredholm determinants. Their quotient is the graded
relative determinant. It is not the ordinary determinant of
\(\cT_0\oplus\cT_1\), which would multiply rather than divide the two
degreewise determinants.
\end{remark}

The two cancellations also have separate ownership. Incidence pair
annihilation removes mixed source labels. The de Rham difference removes a
local holomorphic stability denominator. After both steps, one scalar class
remains per source-derived atom, so the final determinant is again the atom
product.
